\cleardoublepage
\chapterprev{摘要}

光与物质的相互作用在亚波长和亚周期尺度上展现出丰富的物理现象，对其直接测量需要同时具备纳米空间分辨率和亚飞秒时间分辨率。传统光学方法受衍射极限制约，空间分辨率受限；超快电子显微术虽具备原子级空间分辨能力，但时间分辨率受电子脉冲宽度限制，停留在$\sim$\SI{100}{fs}，无法分辨单个光学周期内的动力学过程。

本文发展了基于连续波激光驱动的阿秒电子显微术，在常规透射电子显微镜中实现了$\sim$\SI{800}{as}（阿秒）时间分辨率和$\sim$\SI{50}{nm}空间分辨率的联用。利用有质动力相互作用和时间Talbot效应，将连续波激光的周期性调制转化为电子波函数中的阿秒脉冲结构。通过能量滤波器选择光子诱导能量增益的电子，实现了对纳米结构电磁近场的亚周期时间分辨成像。

本文将阿秒电子显微术应用于三类纳米光子系统的研究：（1）狭缝修饰钨针尖上由圆偏振光激发的手性表面波，实验观测到偏振手性控制表面波传播方向的现象，揭示了自旋-轨道耦合在纳米光子激发中的矢量特性；（2）硅介质纳米共振器中偶极与四极Mie共振模式的时域干涉，首次直接测量了两种模式之间的相位延迟，观测到狭缝内亚光速传播的电磁波前；（3）四狭缝超原子中结构对称性对电磁响应的影响，发现$\sim$\SI{10}{nm}级别的加工偏差足以导致$\pi/3$的相位偏移和电磁能量流的定向偏转。

本文工作建立了阿秒电子显微术作为纳米光子学研究的通用工具，填补了超快科学中高时空分辨测量的关键空白。

\bigskip
\noindent \textbf{关键词}：阿秒电子显微术；超快电子显微镜；近场成像；纳米光子学；高次谐波产生；时间Talbot效应

\cleardoublepage
\chapterprev{Abstract}

The primary interaction between light and matter unfolds on sub-wavelength and sub-cycle dimensions. Direct measurement of these interactions requires simultaneous nanometer spatial resolution and sub-femtosecond temporal resolution. Conventional optical methods are constrained by the diffraction limit, while ultrafast electron microscopy, despite its atomic-scale spatial resolution, has been limited to $\sim$\SI{100}{fs} temporal resolution by electron pulse duration, insufficient for resolving dynamics within a single optical cycle.

This thesis develops attosecond electron microscopy driven by a continuous-wave laser, achieving $\sim$\SI{800}{as} temporal resolution and $\sim$\SI{50}{nm} spatial resolution in a conventional transmission electron microscope. The ponderomotive interaction and temporal Talbot effect convert periodic laser modulation into attosecond pulse structures in the electron wave function. An energy filter selects electrons that have gained energy from photon-induced near-field interactions, enabling sub-cycle time-resolved imaging of electromagnetic near-fields around nanostructures.

The technique is applied to three nanophotonic systems: (1) chiral surface waves on a slit-modified tungsten nanotip, where circular polarization handedness controls wave propagation direction via spin-orbit coupling; (2) temporal interference between dipolar and quadrupolar Mie resonances in a silicon nanoresonator, revealing phase delay between modes and subluminal buried wavefront propagation; (3) symmetry-breaking in a four-slit meta-atom, where $\sim$\SI{10}{nm} fabrication imperfections cause $\pi/3$ phase shifts and redirect electromagnetic energy flow.

This work establishes attosecond electron microscopy as a versatile tool for nanophotonics research, bridging the critical gap in high spatiotemporal resolution measurement.

\bigskip
\noindent \textbf{Keywords}: attosecond electron microscopy; ultrafast electron microscopy; near-field imaging; nanophotonics; high-harmonic generation; temporal Talbot effect
